\documentclass[10pt]{article}

\usepackage[margin=0.75in]{geometry}
\usepackage{amsmath,amsthm,amssymb,bm}
\usepackage{xcolor}
\usepackage{cancel}
\usepackage{graphicx}
\usepackage{changepage}
\usepackage{circuitikz}
\usepackage{pgfplots}
\usepackage{physics}
\usepackage{hyperref}
\usepackage{siunitx}
\usepackage{minted}
\usepackage{fontspec}
\usepackage{relsize}
\usepackage{subfig}
\usepackage{mhchem}
\usepackage{todonotes}
\usepackage{biblatex}
\usepackage{multicol, multirow, booktabs}
\usepackage[breakable]{tcolorbox}
\usepackage[inline]{enumitem}
\patchcmd{\thebibliography}{\section*{\refname}}{}{}{}
\addbibresource{sources.bib}

\theoremstyle{definition}
\newtheorem{problem}{Problem}
\newtheorem{soln}{Solution}

\pgfplotsset{compat=newest}
\usetikzlibrary{lindenmayersystems}
\usetikzlibrary{arrows}
\usetikzlibrary{calc}
\usetikzlibrary{positioning, fit}
\usetikzlibrary{3d, perspective}

\definecolor{incolor}{HTML}{303F9F}
\definecolor{outcolor}{HTML}{D84315}
\definecolor{cellborder}{HTML}{CFCFCF}
\definecolor{cellbackground}{HTML}{F7F7F7}
\newcommand{\ui}{\hat{i}}
\newcommand{\uj}{\hat{j}}
\newcommand{\uk}{\hat{k}}
\newcommand{\ux}{\hat{x}}
\newcommand{\uy}{\hat{y}}
\newcommand{\uz}{\hat{z}}
\newcommand{\uv}[1]{\hat{\bm{{#1}}}}
\newcommand{\pr}[1]{{#1^\prime}}
\newcommand{\justif}[2]{&{#1}&\text{#2}}
\newcommand{\dul}[1]{\underline{\underline{#1}}}
\pgfdeclarelayer{background}  
\pgfsetlayers{background,main}
\AtBeginDocument{\RenewCommandCopy\qty\SI}

\makeatletter
\newcommand{\boxspacing}{\kern\kvtcb@left@rule\kern\kvtcb@boxsep}
\makeatother
\newcommand{\prompt}[4]{
    \ttfamily\llap{{\color{#2}[#3]:\hspace{3pt}#4}}\vspace{-\baselineskip}
}

\newcommand{\thevenin}[2]{
  \begin{center}
    \begin{circuitikz} \draw
      (0,0) -- (2,0) to[battery1, l_=$V_{Th}\eq#1$] (2,2) 
      to[resistor, l_=$R_{Th}\eq#2$] (0,2)
      ;
      \draw [o-] (-.07,2.079);
      \draw [o-] (-.07,0.079);
    \end{circuitikz}
  \end{center}
}

\newcommand{\norton}[2]{
  \begin{center}
    \begin{circuitikz} \draw
      (0,0) -- (3,0) to[american current source, l_=$I_{N}\eq#1$] (3,2) -- (0,2) (2,0)
      to[resistor, l=$R_{N}\eq#2$] (2,2)
      ;
      \draw [o-] (-.07,2.079);
      \draw [o-] (-.07,0.079);
    \end{circuitikz}
  \end{center}
}

\newcommand{\highlight}[1]{\colorbox{yellow}{$\displaystyle #1$}}

\newcommand{\ti}[1]{\widetilde{#1}}

\newfontface{\Kaufmann}{Kaufmann}
\DeclareTextFontCommand{\kf}{\Kaufmann}
\newcommand{\scriptr}{\fontsize{12pt}{12pt}\kf{r}}

\newfontface{\KaufmannB}{Kaufmann Bold}
\DeclareTextFontCommand{\kfb}{\KaufmannB}
\newcommand{\bscriptr}{\fontsize{12pt}{12pt}\kfb{r}}

\newcommand{\bv}[1]{\vec{#1}}

\title{Physics 4605H: Assignment III}
\author{Jeremy Favro (0805980) \\ Trent University, Peterborough, ON, Canada}
\date{\today}

\begin{document}
\maketitle

% PROBLEM 0
\setcounter{problem}{0}
\setcounter{soln}{0}
\begin{problem}
Consider a spin-$1/2$ object in a uniform magnetic field $\bv{B}=B_0\uz$. Ignore the chemical
shift, such that
$$\hat{H}=-\bv{\mu}\cdot\bv{B}=-\frac{\gamma \hbar B_0}{2}\hat{\sigma}_z.$$
Let the initial state of the particle be
$$\ket{\psi(t=0)}=\frac{1}{\sqrt{2}}\ket{0}+\frac{1}{\sqrt{2}}\ket{1}$$
Now consider the time dependence. Recall that $\ket{0}$ and $\ket{1}$ are eigenstates of the Hamiltonian.
\begin{enumerate}[label=(\alph*)]
  \setcounter{enumi}{5}
  \item What is the wavefunction at a later time $t$?
  \item What is the expectation value of $S_x$ as a function of time?
  \item What is the expectation value of $S_y$ as a function of time?
  \item What is the expectation value of $S_z$ as a function of time?
  \item At time $t = 0$, as shown in (c)-(e), the expectation value of $\bv{S}$ is entirely along $x$. What
        is the (first) time at which it will be entirely along $y$?
  \item In what direction does it rotate?
\end{enumerate}
\end{problem}
\begin{soln}
  \begin{enumerate}[label=(\alph*)]
    \setcounter{enumi}{5}
    \item $$\ket{\psi(x,t)}=\frac{1}{\sqrt{2}}\ket{0}e^{i\frac{\gamma B_0}{2}}+\frac{1}{\sqrt{2}}\ket{1}e^{-i\frac{\gamma B_0}{2}}$$
          where we've used the eigenvalues $\mp\frac{\gamma \hbar B_0}{2}$ of the Hamiltonian (the energies), corresponding to $\ket{0}$ and $\ket{1}$ respectively,s in $e^{-iEt/\hbar}$.
    \item
          \begin{align*}
            \langle S_x\rangle & =\bra{\psi}\hat{S}_x\ket{\psi}=(\frac{1}{\sqrt{2}}\bra{0}e^{-i\frac{\gamma B_0t}{2}}+\frac{1}{\sqrt{2}}\bra{1}e^{i\frac{\gamma B_0t}{2}})
            \hat{S}_x(\frac{1}{\sqrt{2}}\ket{0}e^{i\frac{\gamma B_0t}{2}}+\frac{1}{\sqrt{2}}\ket{1}e^{-i\frac{\gamma B_0t}{2}})                                                                                                                                     \\
                               & =\frac{\hbar}{4}(\bra{0}e^{-i\frac{\gamma B_0t}{2}}+\bra{1}e^{i\frac{\gamma B_0t}{2}})
            (\ket{1}e^{i\frac{\gamma B_0t}{2}}+\ket{0}e^{-i\frac{\gamma B_0t}{2}})                                                                                                                                                                                  \\
                               & =\frac{\hbar}{4}(\bra{0}(\ket{1}e^{i\frac{\gamma B_0t}{2}}+\ket{0}e^{-i\frac{\gamma B_0t}{2}})e^{-i\frac{\gamma B_0t}{2}}+\bra{1}(\ket{1}e^{i\frac{\gamma B_0t}{2}}+\ket{0}e^{-i\frac{\gamma B_0t}{2}})e^{i\frac{\gamma B_0t}{2}}) \\
                               & =\frac{\hbar}{4}((e^{-i\frac{\gamma B_0t}{2}})e^{-i\frac{\gamma B_0t}{2}}+(e^{i\frac{\gamma B_0t}{2}})e^{i\frac{\gamma B_0t}{2}})                                                                                                  \\
                               & =\frac{\hbar}{4}(e^{-i\gamma B_0t}+e^{i\gamma B_0t})=\frac{\hbar}{2}\cos(\gamma B_0t)
          \end{align*}
    \item  \begin{align*}
            \langle S_y\rangle & =\bra{\psi}\hat{S}_y\ket{\psi}=(\frac{1}{\sqrt{2}}\bra{0}e^{-i\frac{\gamma B_0t}{2}}+\frac{1}{\sqrt{2}}\bra{1}e^{i\frac{\gamma B_0t}{2}})
            \hat{S}_y(\frac{1}{\sqrt{2}}\ket{0}e^{i\frac{\gamma B_0t}{2}}+\frac{1}{\sqrt{2}}\ket{1}e^{-i\frac{\gamma B_0t}{2}})                                                                                                                                           \\
                               & =\frac{\hbar}{4}(\bra{0}e^{-i\frac{\gamma B_0t}{2}}+\bra{1}e^{i\frac{\gamma B_0t}{2}})
            (i\ket{1}e^{i\frac{\gamma B_0t}{2}}-i\ket{0}e^{-i\frac{\gamma B_0t}{2}})                                                                                                                                                                                      \\
                               & =\frac{\hbar}{4}(\bra{0}(i\ket{1}e^{i\frac{\gamma B_0t}{2}}-i\ket{0}e^{-i\frac{\gamma B_0t}{2}}) e^{-i\frac{\gamma B_0t}{2}}+\bra{1}(i\ket{1}e^{i\frac{\gamma B_0t}{2}}-i\ket{0}e^{-i\frac{\gamma B_0t}{2}}) e^{i\frac{\gamma B_0t}{2}}) \\
                               & =\frac{\hbar}{4}((-ie^{-i\frac{\gamma B_0t}{2}}) e^{-i\frac{\gamma B_0t}{2}}+(ie^{i\frac{\gamma B_0t}{2}}) e^{i\frac{\gamma B_0t}{2}})                                                                                                   \\
                               & =\frac{i\hbar}{4}(e^{i\gamma B_0t}-e^{-i\gamma B_0t})=\frac{\hbar}{2}\sin(\gamma B_0t)
          \end{align*}
    \item \begin{align*}
            \langle S_z\rangle & =\bra{\psi}\hat{S}_z\ket{\psi}=(\frac{1}{\sqrt{2}}\bra{0}e^{-i\frac{\gamma B_0t}{2}}+\frac{1}{\sqrt{2}}\bra{1}e^{i\frac{\gamma B_0t}{2}})
            \hat{S}_z(\frac{1}{\sqrt{2}}\ket{0}e^{i\frac{\gamma B_0t}{2}}+\frac{1}{\sqrt{2}}\ket{1}e^{-i\frac{\gamma B_0t}{2}})                                                                                                                                     \\
                               & =\frac{\hbar}{4}(\bra{0}e^{-i\frac{\gamma B_0t}{2}}+\bra{1}e^{i\frac{\gamma B_0t}{2}})
            (\ket{0}e^{i\frac{\gamma B_0t}{2}}-\ket{1}e^{-i\frac{\gamma B_0t}{2}})                                                                                                                                                                                  \\
                               & =\frac{\hbar}{4}(\bra{0}(\ket{0}e^{i\frac{\gamma B_0t}{2}}-\ket{1}e^{-i\frac{\gamma B_0t}{2}})e^{-i\frac{\gamma B_0t}{2}}+\bra{1}(\ket{0}e^{i\frac{\gamma B_0t}{2}}-\ket{1}e^{-i\frac{\gamma B_0t}{2}})e^{i\frac{\gamma B_0t}{2}}) \\
                               & =\frac{\hbar}{4}(e^{i\frac{\gamma B_0t}{2}}e^{-i\frac{\gamma B_0t}{2}}-e^{-i\frac{\gamma B_0t}{2}}e^{i\frac{\gamma B_0t}{2}})                                                                                                      \\
                               & =\frac{\hbar}{4}(1-1)=0
          \end{align*}
    \item As shown in (g) \& (h) the expectation value is governed by cosine and sine. Hence when the argument of the $\frac{\hbar}{2}\sin(\gamma B_0t)$ term is
          $\gamma B_0t =\frac{\pi}{2}\implies t=\frac{\pi}{2B_0\gamma}$ $\langle S_x\rangle=0$ and $\langle S_y \rangle\frac{\hbar}{2}$.
    \item From $x\to y$, or in spherical coordinates, the $+\uv{\theta}$ direction.
  \end{enumerate}

\end{soln}

% PROBLEM 1
\begin{problem}
Consider a single-qubit operator which, in the basis $\ket{0},\ket{1}$, is
$$\hat{A}=\frac{1}{\sqrt{2}}\begin{pmatrix}
    1 & 1  \\
    1 & -1
  \end{pmatrix}$$
\begin{enumerate}[label=(\alph*)]
  \item Write this operator in terms of Pauli operators.
  \item Write this operator in Dirac notation.
  \item What is $\hat{A}\ket{0}$?
  \item What is $\hat{A}^2\ket{0}$?
\end{enumerate}
\end{problem}
\begin{soln}~
  \begin{enumerate}[label=(\alph*)]
    \item $$\hat{A}=\frac{1}{\sqrt{2}}\left[ \dul{\sigma}_x+\dul{\sigma}_z\right]$$
    \item We can expand the matrix as
          $$\begin{pmatrix}
              1 & 1  \\
              1 & -1
            \end{pmatrix}=\begin{pmatrix}
              1 & 0 \\
              0 & 0
            \end{pmatrix}+\begin{pmatrix}
              0 & 1 \\
              0 & 0
            \end{pmatrix}+\begin{pmatrix}
              0 & 0 \\
              1 & 0
            \end{pmatrix}-\begin{pmatrix}
              0 & 0 \\
              0 & 1
            \end{pmatrix}$$
          which can clearly be written as
          $$\ket{0}\bra{0}+\ket{0}\bra{1}+\ket{1}\bra{0}-\ket{1}\bra{1}=\ket{0}(\bra{0}+\bra{1})+\ket{1}(\bra{0}-\bra{1})$$
    \item This is easiest done using our expression in Dirac notation,
          \begin{align*}
            \hat{A}\ket{0} & =\frac{1}{\sqrt{2}}(\ket{0}(\bra{0}+\bra{1})+\ket{1}(\bra{0}-\bra{1}))\ket{0}        \\
                           & =\frac{1}{\sqrt{2}}(\ket{0}(\bra{0}+\bra{1})\ket{0}+\ket{1}(\bra{0}-\bra{1})\ket{0}) \\
                           & =\frac{1}{\sqrt{2}}(\ket{0}(1)+\ket{1}(1))=\frac{1}{\sqrt{2}}\begin{pmatrix}
                                                                                            1 \\1
                                                                                          \end{pmatrix}          \\
          \end{align*}
    \item We can just apply $\hat{A}$ to our previous result to obtain
          \begin{align*}
            \hat{A}^2\ket{0} & =\hat{A}\frac{1}{\sqrt{2}}(\ket{0}+\ket{1}) \\
                             & =\frac{1}{2}\begin{pmatrix}
                                             1 & 1  \\
                                             1 & -1
                                           \end{pmatrix}\begin{pmatrix}
                                                          1 \\1
                                                        \end{pmatrix}     \\
                             & =\frac{1}{2}\begin{pmatrix}
                                             2 \\0
                                           \end{pmatrix}=\ket{0}
          \end{align*}
  \end{enumerate}
\end{soln}

% PROBLEM 2
\begin{problem}
The following states are known as Bell states:
\begin{align*}
  \ket{B_{00}}=\left(\ket{00}+\ket{11}\right)/\sqrt{2};\quad & \ket{B_{01}}=\left(\ket{00}-\ket{11}\right)/\sqrt{2} \\
  \ket{B_{10}}=\left(\ket{01}+\ket{10}\right)/\sqrt{2};\quad & \ket{B_{11}}=\left(\ket{01}-\ket{10}\right)/\sqrt{2}
\end{align*}
You may assume that the set of states $\ket{00}$, $\ket{01}$, $\ket{10}$, $\ket{11}$ is complete and orthonormal.
\begin{enumerate}[label=(\alph*)]
  \item Show that $\ket{B_{00}}$ is normalized.
  \item Show that $\ket{B_{00}}$ and $\ket{B_{01}}$ are orthogonal.
  \item Show that these states form a complete set for a 2-qubit Hilbert space.
\end{enumerate}
\end{problem}
\begin{soln}
  \begin{enumerate}[label=(\alph*)]
    \item \begin{align*}
            \bra{B_{00}}\ket{B_{00}} & =\left(\bra{00}+\bra{11}\right)/\sqrt{2}\left(\ket{00}+\ket{11}\right)/\sqrt{2}                                        \\
                                     & =\frac{1}{2}\left(\left(\bra{00}+\bra{11}\right)\ket{00}+\left(\bra{00}+\bra{11}\right)\ket{11}\right)                 \\
                                     & =\frac{1}{2}\left(\left(\bra{00}\ket{00}+\bra{11}\ket{00}\right)+\left(\bra{00}\ket{11}+\bra{11}\ket{11}\right)\right) \\
                                     & =\frac{1}{2}\left(1+1\right)=1
          \end{align*}
    \item \begin{align*}
            \bra{B_{00}}\ket{B_{01}} & =\frac{1}{2}\left(\bra{00}+\bra{11}\right)\left(\ket{00}-\ket{11}\right)                                               \\
                                     & =\frac{1}{2}\left(\left(\bra{00}+\bra{11}\right)\ket{00}-\left(\bra{00}+\bra{11}\right)\ket{11}\right)                 \\
                                     & =\frac{1}{2}\left(\left(\bra{00}\ket{00}+\bra{11}\ket{00}\right)-\left(\bra{00}\ket{11}+\bra{11}\ket{11}\right)\right) \\
                                     & =\frac{1}{2}\left(\left(1+0\right)-\left(0+1\right)\right)=0
          \end{align*}
    \item Here we want to show that the completeness relation holds, i.e. that
          $$\ket{B_{00}}\bra{B_{00}}+\ket{B_{01}}\bra{B_{01}}+\ket{B_{10}}\bra{B_{10}}+\ket{B_{11}}\bra{B_{11}}$$
          is the identity. We can evaluate each of these pairs,
          $$\ket{B_{00}}\bra{B_{00}}=\begin{pmatrix}
              1 \\0\\0\\0
            \end{pmatrix}\begin{pmatrix}
              1 & 0 & 0 & 0
            \end{pmatrix}=\begin{pmatrix}
              1 & 0 & 0 & 0 \\
              0 & 0 & 0 & 0 \\
              0 & 0 & 0 & 0 \\
              0 & 0 & 0 & 0
            \end{pmatrix},$$
          $$\ket{B_{01}}\bra{B_{01}}=\begin{pmatrix}
              0 \\1\\0\\0
            \end{pmatrix}\begin{pmatrix}
              0 & 1 & 0 & 0
            \end{pmatrix}=\begin{pmatrix}
              0 & 0 & 0 & 0 \\
              0 & 1 & 0 & 0 \\
              0 & 0 & 0 & 0 \\
              0 & 0 & 0 & 0
            \end{pmatrix},$$
          $$\ket{B_{00}}\bra{B_{00}}=\begin{pmatrix}
              0 \\0\\1\\0
            \end{pmatrix}\begin{pmatrix}
              0 & 0 & 1 & 0
            \end{pmatrix}=\begin{pmatrix}
              0 & 0 & 0 & 0 \\
              0 & 0 & 0 & 0 \\
              0 & 0 & 1 & 0 \\
              0 & 0 & 0 & 0
            \end{pmatrix},$$
          $$\ket{B_{00}}\bra{B_{00}}=\begin{pmatrix}
              0 \\0\\0\\1
            \end{pmatrix}\begin{pmatrix}
              0 & 0 & 0 & 1
            \end{pmatrix}=\begin{pmatrix}
              0 & 0 & 0 & 0 \\
              0 & 0 & 0 & 0 \\
              0 & 0 & 0 & 0 \\
              0 & 0 & 0 & 1
            \end{pmatrix}.$$
            It is evident that these sum to the identity.
  \end{enumerate}
\end{soln}
\end{document}